\begin{ex}
Hình tam giác có bao nhiêu cạnh?
\shortans{3}
\loigiai{
Hình tam giác có đúng 3 cạnh.
\textbf{Đáp số:} $3$.
}
\end{ex}

\begin{ex}
Hình chữ nhật có bao nhiêu góc vuông?
\shortans{4}
\loigiai{
Hình chữ nhật có đúng 4 góc vuông.
\textbf{Đáp số:} $4$.
}
\end{ex}

\begin{ex}
Khối lập phương có bao nhiêu mặt?
\shortans{6}
\loigiai{
Khối lập phương có đúng 6 mặt phẳng đều là hình vuông.
\textbf{Đáp số:} $6$.
}
\end{ex}

\begin{ex}
Khối hộp chữ nhật có bao nhiêu đỉnh?
\shortans{8}
\loigiai{
Khối hộp chữ nhật có đúng 8 đỉnh.
\textbf{Đáp số:} $8$.
}
\end{ex}

\begin{ex}
Đồng hồ có kim ngắn chỉ số 8, kim dài chỉ số 12. Hỏi đồng hồ chỉ mấy giờ?
\shortans{8}
\loigiai{
Đồng hồ chỉ đúng 8 giờ.
\textbf{Đáp số:} $8$.
}
\end{ex}

\begin{ex}
Một tuần lễ có bao nhiêu ngày?
\shortans{7}
\loigiai{
Một tuần lễ có đúng 7 ngày.
\textbf{Đáp số:} $7$.
}
\end{ex}

\begin{ex}
Tính: $35\text{ cm} + 24\text{ cm}$. Kết quả là bao nhiêu xăng-ti-mét?
\shortans{59}
\loigiai{
Ta có: $35 + 24 = 59$ (cm).
\textbf{Đáp số:} $59$.
}
\end{ex}

\begin{ex}
Tính: $78\text{ cm} - 35\text{ cm}$. Kết quả là bao nhiêu xăng-ti-mét?
\shortans{43}
\loigiai{
Ta có: $78 - 35 = 43$ (cm).
\textbf{Đáp số:} $43$.
}
\end{ex}

\begin{ex}
Một sợi dây dài 65 cm. Bạn Mai cắt bớt 22 cm. Hỏi sợi dây còn lại dài bao nhiêu xăng-ti-mét?
\shortans{43}
\loigiai{
Chiều dài còn lại là: $65 - 22 = 43$ (cm).
\textbf{Đáp số:} $43$.
}
\end{ex}

\begin{ex}
Đoạn thẳng AB dài 40 cm, đoạn thẳng CD dài 50 cm. Cả hai đoạn thẳng dài bao nhiêu xăng-ti-mét?
\shortans{90}
\loigiai{
Tổng độ dài hai đoạn là: $40 + 50 = 90$ (cm).
\textbf{Đáp số:} $90$.
}
\end{ex}

\begin{ex}
Nếu hôm nay là Thứ Tư ngày 12 thì Thứ Tư tuần sau là ngày bao nhiêu?
\shortans{19}
\loigiai{
Thứ Tư tuần sau là: $12 + 7 = 19$.
\textbf{Đáp số:} $19$.
}
\end{ex}

\begin{ex}
Nếu hôm nay là Thứ Bảy ngày 25 thì Thứ Bảy tuần trước là ngày bao nhiêu?
\shortans{18}
\loigiai{
Thứ Bảy tuần trước là: $25 - 7 = 18$.
\textbf{Đáp số:} $18$.
}
\end{ex}

\begin{ex}
Một ngày có bao nhiêu giờ?
\shortans{24}
\loigiai{
Một ngày có đúng 24 giờ.
\textbf{Đáp số:} $24$.
}
\end{ex}

\begin{ex}
Cần ít nhất bao nhiêu que tính để xếp thành 2 hình tam giác riêng biệt?
\shortans{6}
\loigiai{
Mỗi hình tam giác cần 3 que: $3 + 3 = 6$ (que tính).
\textbf{Đáp số:} $6$.
}
\end{ex}

\begin{ex}
Cần ít nhất bao nhiêu que tính để xếp thành 2 hình vuông riêng biệt?
\shortans{8}
\loigiai{
Mỗi hình vuông cần 4 que: $4 + 4 = 8$ (que tính).
\textbf{Đáp số:} $8$.
}
\end{ex}

\begin{ex}
Hình tròn có bao nhiêu cạnh thẳng?
\shortans{0}
\loigiai{
Hình tròn không có cạnh thẳng nào (có 0 cạnh thẳng).
\textbf{Đáp số:} $0$.
}
\end{ex}

\begin{ex}
Một hình chữ nhật dài 25 cm, rộng 14 cm. Chiều dài hơn chiều rộng bao nhiêu xăng-ti-mét?
\shortans{11}
\loigiai{
Chiều dài hơn chiều rộng là: $25 - 14 = 11$ (cm).
\textbf{Đáp số:} $11$.
}
\end{ex}

\begin{ex}
Bác thợ mộc có 8 khối gỗ hình lập phương. Bác dùng 5 khối gỗ đóng giá sách. Bác còn lại bao nhiêu khối gỗ?
\shortans{3}
\loigiai{
Số khối gỗ còn lại là: $8 - 5 = 3$ (khối).
\textbf{Đáp số:} $3$.
}
\end{ex}

\begin{ex}
Đồng hồ chỉ 2 giờ chiều. Sau 4 giờ nữa đồng hồ chỉ mấy giờ chiều?
\shortans{6}
\loigiai{
Thời điểm đó là: $2 + 4 = 6$ (giờ chiều).
\textbf{Đáp số:} $6$.
}
\end{ex}

\begin{ex}
Đồng hồ chỉ 11 giờ trưa. Trước đó 3 giờ là mấy giờ sáng?
\shortans{8}
\loigiai{
Thời điểm đó là: $11 - 3 = 8$ (giờ sáng).
\textbf{Đáp số:} $8$.
}
\end{ex}

\begin{ex}
Một tuần lễ em đi học 5 ngày. Hỏi trong 4 tuần lễ em đi học bao nhiêu ngày?
\shortans{20}
\loigiai{
Tổng số ngày đi học là: $5 + 5 + 5 + 5 = 20$ (ngày).
\textbf{Đáp số:} $20$.
}
\end{ex}

\begin{ex}
Một băng giấy dài 80 cm, cắt thành hai đoạn bằng nhau. Mỗi đoạn dài bao nhiêu xăng-ti-mét?
\shortans{40}
\loigiai{
Mỗi đoạn dài $40$ cm (vì $40 + 40 = 80$).
\textbf{Đáp số:} $40$.
}
\end{ex}

\begin{ex}
Tổng số mặt của hai khối lập phương là bao nhiêu mặt?
\shortans{12}
\loigiai{
Mỗi khối có 6 mặt: $6 + 6 = 12$ (mặt).
\textbf{Đáp số:} $12$.
}
\end{ex}

\begin{ex}
Tổng số đỉnh của một hình tam giác và một khối hộp chữ nhật là bao nhiêu?
\shortans{11}
\loigiai{
Hình tam giác có 3 đỉnh, khối hộp chữ nhật có 8 đỉnh: $3 + 8 = 11$ (đỉnh).
\textbf{Đáp số:} $11$.
}
\end{ex}

\begin{ex}
Một thước kẻ dài 20 cm, đặt que tính dài 8 cm lên thước từ vạch số 0. Đầu kia que tính chỉ vạch số mấy?
\shortans{8}
\loigiai{
Đầu kia chỉ đúng vạch số $8$.
\textbf{Đáp số:} $8$.
}
\end{ex}
